%%% CSE 4316 Project Charter - WORKED EXAMPLE (Autonomous Ground Robot)
%%% Structure follows ISO/IEC/IEEE 16326:2019. See ../../project charter latex/ for the blank template.
\documentclass[11pt,letterpaper]{article}
\usepackage[margin=1.0in]{geometry}
\usepackage[T1]{fontenc}
\usepackage[bitstream-charter]{mathdesign}
\usepackage[latin1]{inputenc}
\usepackage{amsmath}\usepackage{xcolor}\usepackage{cite}\usepackage{hyphenat}
\usepackage{graphicx}\usepackage{float}\usepackage{subfigure}\usepackage{sectsty}
\usepackage[compact]{titlesec}\usepackage[tablegrid]{vhistory}\usepackage{longtable}
\allsectionsfont{\color{accentcolor}\scshape\selectfont}
\definecolor{accentcolor}{rgb}{0.0,0.0,0.5}
\newcommand{\teamname}{Team Trailblazer}
\newcommand{\productname}{TerraScout Autonomous Ground Rover}
\newcommand{\coursename}{CSE 4316: Senior Design I}
\newcommand{\semester}{Fall 2026}
\newcommand{\docname}{Project Charter}
\newcommand{\department}{Department of Computer Science \& Engineering}
\newcommand{\university}{The University of Texas at Arlington}
\newcommand{\authors}{Alan Turing \\ Grace Hopper \\ John Von Neumann \\ Ada Lovelace \\ Charles Babbage}
\usepackage{fancyhdr}\pagestyle{fancy}\usepackage{lastpage}
\lhead{}\chead{}\rhead{}\lfoot{\footnotesize \teamname \ - \semester}\cfoot{}
\rfoot{\footnotesize page \thepage\ of \pageref{LastPage}}
\renewcommand{\headrulewidth}{0.0pt}\renewcommand{\footrulewidth}{0.4pt}
\usepackage[runin]{abstract}\abslabeldelim{\quad}\setlength{\abstitleskip}{-10pt}
\renewcommand{\abstractname}{}\renewcommand{\abstracttextfont}{\color{accentcolor} \small \slshape}

\begin{document}
{\centering \huge \color{accentcolor} \sc \textbf{\department \\ \university} \par}
\vspace{1 in}
{\centering \huge \color{accentcolor} \sc \textbf{\docname \\ \coursename \\ \semester} \par}
\vspace{0.5 in}
\begin{figure}[h!]\centering\includegraphics[width=0.60\textwidth]{images/test_image}\end{figure}
\vspace{0.5 in}
{\centering \huge \color{accentcolor} \sc \textbf{\teamname \\ \productname} \par}
\vspace{0.5 in}
{\centering \large \sc \textbf{\authors} \par}
\newpage

\begin{versionhistory}
  \vhEntry{0.1}{09.02.2026}{GH}{document creation}
  \vhEntry{0.2}{09.06.2026}{AT|GH}{complete draft}
  \vhEntry{1.0}{09.15.2026}{AT|GH|CB}{official release}
\end{versionhistory}
\newpage
\tableofcontents
\newpage

\section{Problem Statement}
Utility-scale solar farms span hundreds of acres and thousands of photovoltaic (PV)
panels. Undetected panel faults --- cracked cells, soiling, and thermal hotspots ---
silently reduce energy yield and can become fire hazards. Today these sites are
inspected manually on foot or with piloted drones, which is slow, labor-intensive,
weather-limited, and inconsistent between inspectors. The result is that faults
often go unaddressed for weeks. What is needed is a way to survey a PV field
frequently, consistently, and without tying up skilled staff, so that faults are
found and fixed while the lost energy is still small.

\section{Methodology}
Team Trailblazer will build \productname, a small four-wheeled autonomous ground
rover that patrols the drivable lanes between PV rows. Following a set of GPS
waypoints, the rover captures visual and thermal imagery of the panels, flags
likely faults on-board, and produces a GPS-tagged inspection report at the end of
each run. When it detects a low battery or an obstacle it cannot pass, it returns
to its charging dock autonomously. In short: we are building a robot that drives
itself around a solar farm and tells the operator which panels need attention.

\section{Value Proposition}
For the sponsor (a regional utility's O\&M group), the rover converts a costly,
irregular manual task into a repeatable automated one, catching yield-robbing
faults earlier and reducing truck-rolls to healthy panels. For the university and
instructors, the project exercises the full systems-engineering life cycle across
embedded hardware, autonomy software, and computer vision, producing a
demonstrable prototype and a reusable mobile-robot platform for future teams.

\section{Development Milestones}
\begin{itemize}
  \item Project Charter first draft - September 2026
  \item System Requirements Specification - October 2026
  \item Architectural Design Specification - November 2026
  \item Demonstration of tele-operated drive \& sensor capture - December 2026
  \item Detailed Design Specification - February 2027
  \item Demonstration of waypoint navigation - March 2027
  \item Demonstration of on-board fault detection - March 2027
  \item CoE Innovation Day poster presentation - April 2027
  \item Demonstration of autonomous dock/return - April 2027
  \item Final Project Demonstration (full autonomous patrol) - May 2027
\end{itemize}
\newpage

\section{Background}
The sponsor operates several solar sites and currently contracts periodic drone
thermography, typically once or twice a year. Between contracted flights the panels
go uninspected. The sponsor approached the CSE Senior Design program seeking a
lower-cost, higher-frequency ground-based alternative that their own field staff
could operate without a licensed drone pilot. The team has no prior relationship
with the sponsor; requirements were gathered through two scoping interviews.

\section{Related Work}
Aerial PV inspection with thermographic drones is well established
\cite{Rubin2012}, but requires certified pilots and favorable weather, and does
not run unattended. Commercial autonomous ground robots for solar exist but are
priced for large operators and are closed platforms. Academic mobile-robot
navigation using ROS, GPS-RTK, and visual fault detection is mature enough to
adapt. \productname\ targets the gap: an affordable, operator-runnable ground
platform for frequent inspection. (Add at least five fully-documented IEEE-format
references in the final version.)

\section{System Overview}
At a high level the rover comprises a sensing front end (RGB + thermal cameras,
GPS, IMU, ultrasonic bumpers), an on-board compute module running the autonomy and
detection software, and a mobility base (brushless drive motors, battery, motor
controllers). The operator interacts through a laptop ground-station web page over
Wi-Fi to start runs and review reports; a charging dock provides autonomous
recharge. See the Architectural Design Specification for the layered decomposition.

\section{Project Management Approach}
The team follows an iterative Scrum model \cite{Rubin2012, SchwaberSutherland2020}
consistent with ISO/IEC/IEEE 16326:2019 \cite{ISO16326} within the life-cycle
frameworks of ISO/IEC/IEEE 12207:2017 \cite{ISO12207} and 15288:2023
\cite{ISO15288}. Two-week sprints run across CSE 4316 (requirements/design) and
CSE 4317 (build/test). Work is estimated in task-hours and tracked on a shared
board; progress is measured by sprint burndown and requirement/test coverage. A
burndown trending above the ideal line for three consecutive days triggers scope
re-negotiation with the Product Owner. Requirement changes are approved by the
Product Owner in consultation with the sponsor; all work products are version-
controlled in Git with tagged document releases.

\section{Roles \& Responsibilities}
Sponsor point of contact: utility O\&M lead. Product Owner: Grace Hopper.
Scrum Master: Alan Turing (rotating each semester). Areas of responsibility ---
Perception/CV: Ada Lovelace; Navigation/autonomy: John Von Neumann; Embedded
platform/power: Charles Babbage; Ground station/reporting: Grace Hopper;
Integration/test: Alan Turing.

\section{Cost Proposal}
Major expenses are the compute module, the thermal camera (the single largest line
item), drive motors and controllers, battery, and chassis materials. Funding is
the default CSE department stipend supplemented by a sponsor hardware contribution;
the thermal camera may be loaned by the sponsor to stay within budget.

\section{Facilities \& Equipment}
Development uses the Senior Design lab (soldering, bench power, 3D printing). Field
testing uses a university parking area mocked up with panel stand-ins, and, later,
a supervised visit to a sponsor solar site. RTK-GPS base station and a fenced
outdoor test area are required by the third sprint.

\section{Assumptions}
\begin{itemize}
  \item A fenced outdoor testing area will be available by the 3rd sprint.
  \item The sponsor will loan a compatible thermal camera by the 4th sprint.
  \item Panel rows provide a drivable, reasonably level lane ($\leq$10\% grade).
  \item RTK-GPS correction will be available at the test site.
  \item The ground-station laptop and rover share a private Wi-Fi network.
\end{itemize}

\section{Constraints}
\begin{itemize}
  \item Final prototype demonstration must be completed by May 1st, 2027.
  \item Total development cost must not exceed \$800 excluding loaned equipment.
  \item The rover must weigh under 25\,kg for safe two-person handling.
  \item Field testing only during daylight and dry weather.
  \item All captured imagery of sponsor sites must be sponsor-approved before public release.
\end{itemize}

\section{Risks}
\begin{table}[h]
\resizebox{\textwidth}{!}{
\begin{tabular}{|l|l|l|l|}
\hline
\textbf{Risk description} & \textbf{Probability} & \textbf{Loss (days)} & \textbf{Exposure (days)} \\ \hline
Thermal camera loan delayed & 0.40 & 20 & 8.0 \\ \hline
Outdoor RTK test area unavailable & 0.30 & 14 & 4.2 \\ \hline
Autonomy integration harder than estimated & 0.50 & 10 & 5.0 \\ \hline
Motor controller lead time (overseas) & 0.20 & 18 & 3.6 \\ \hline
Weather blocks field test windows & 0.35 & 8 & 2.8 \\ \hline
\end{tabular}}
\caption{Highest-exposure project risks}
\end{table}

\section{Documentation \& Reporting}
Documentation deliverables (Charter, SRS, ADS, DDS, Test Plan) are maintained in
Git and released at the ends of their planning sprints. Recurring artifacts ---
product/sprint backlogs, sprint goals, burndown charts, individual sprint reports,
retrospectives, and engineering notebooks --- are produced each sprint. Closeout
materials include the rover prototype, a project poster, a public web page, a demo
video, documented source code (MIT license), wiring schematics, chassis CAD (STEP),
deployment scripts for the ground station, and an operator manual.

\bibliographystyle{reference/IEEEtran_custom}
\bibliography{reference/refs}{}
\end{document}
